% =====================================================================
%  28.3  Unrooted isomorphism: the centre as canonical root, the
%  two-centre (centre-edge) pairing test, and the chapter's SECOND
%  solvedproblem: CF 1800G "Symmetree" (live-verified: rating 2200,
%  L4, tags dfs/hashing/trees/constructive; statement = rooted at 1,
%  mirror-pairing definition, t<=1e4, n<=2e5, sum n<=2e5; six official
%  samples replayed here -> YES NO YES NO NO YES).
%  Exports: sec:c22-centre  def:c22-canroot  thm:c22-centreiso
%  thm:c22-sym  lst:c22-unrooted  lst:c22-symmetree  prob:c22-symmetree
%  fig:c22-pair  fig:c22-symmetree
%  The centre/diameter COMPUTATION is Ch.12 (\chref{c09}) and is quoted
%  as a black box, never re-derived.  CF 1182/D "Complete Mirror" is
%  named in prose ONLY (its URL belongs to ch20_advanced_tree_dp).
% =====================================================================
\section{Unrooted Isomorphism: Root at the Centre}\label{sec:c22-centre}

\begin{previouslybox}[28.1--28.2 can compare shapes -- but only shapes with a root]
Two tools so far: interned canonical labels (\(\secref{c22-ahu}\)) and randomized
subtree hashes (\(\secref{c22-hash}\)).  Both answer a \emph{rooted} question --
\(\subtree(u)\cong\subtree(w)\) -- and both were built on the recursion order that
a root supplies.  Remove the root from the input and the same question becomes
sharper: \(R\) and its twin \(R'\) (\(\figref{c22-twin}\)) are literally the same
eight edges wearing different names, yet the rooted verdict at vertex \(1\) said NO
(\(4\neq7\)).  The fix is not a cleverer compare; it is to stop inventing a root
and use the one the tree \emph{already has}: the centre, which Ch.~12 taught you
to find from the diameter.  Pairing shapes is the other half of this section, and
CF 1800G is the payoff.
\end{previouslybox}

\begin{topicintro}
An isomorphism cannot move the centre: the centre is \emph{defined} out of the
tree by a rule that mentions no labels (minimize the maximum distance), so every
isomorphism maps centre to centre.  Root both trees there and the rooted
machinery of 28.1 applies verbatim -- equal root names now mean equal
\emph{unrooted} trees, and the wrong-root lie of \(R\) versus \(R'\) disappears
(\(\thmref{c22-centreiso}\)).  One wrinkle: the centre is sometimes an
\emph{edge}, and an edge has two endpoints that the tree cannot tell apart.  The
answer is to keep both rooted forms and compare them as a pair -- and that is
exactly the same move CF 1800G needs at every internal vertex, where the children
must be paired into mirror images with (at most) one self-similar child in the
middle (\(\thmref{c22-sym}\)).

\textbf{Prerequisites.} Chapter~12 (\chref{c09}) -- diameter, eccentricity and
the theorem that a tree has one centre vertex or one centre edge; used as a black
box.  \(\secref{c22-ahu}\) for the labeller, \(\secref{c22-hash}\) if you prefer
hashes as the pairing key.

\textbf{After this section you can}
\begin{itemize}[nosep]
  \item state why the centre is the only label-free canonical root, and root both trees there before comparing; \hfill\lv{3}
  \item handle the two-centre case with the paired-form test instead of picking an endpoint; \hfill\lv{4}
  \item decide unrooted isomorphism of two trees in \(\bigO{n\log n}\) with one shared id table; \hfill\lv{4}
  \item decide mirror-symmetry of a rooted tree with the "at most one odd child class" test (solved CF 1800G); \hfill\lv{4}
  \item explain why equal subtree \emph{sizes} are not equal \emph{shapes} -- the classic WA of 1800G. \hfill\lv{3}
\end{itemize}
\end{topicintro}

% ---------------------------------------------------------------
\subsection{The centre is the only honest root}
% ---------------------------------------------------------------

\begin{definition}{Canonical form \(\Phi(T)\) of an unrooted tree}{c22-canroot}
Let \(\mathrm{centre}(T)\) be the centre (Ch.~12): one vertex, or the two
endpoints of one edge.  If \(\mathrm{centre}(T)=\{c\}\), set
\(\Phi(T)=\mathrm{cl}(T\text{ rooted at }c)\).  If
\(\mathrm{centre}(T)=\{a,b\}\) (the centre edge), set
\[
   \Phi(T)\;=\;\bigl(\min(x,y),\ \max(x,y)\bigr),\qquad
   x=\mathrm{cl}(T\text{ rooted at }a),\quad y=\mathrm{cl}(T\text{ rooted at }b),
\]
the \emph{sorted pair} of the two endpoint-rooted names.  All labellings of all
trees being compared run against ONE shared \(\mathit{tid}\) table
(\(\secref{c22-ahu}\), pitfall~1).
\end{definition}

\begin{theorem}{\(T\cong T'\) iff \(\Phi(T)=\Phi(T')\)}{c22-centreiso}
Every isomorphism preserves distances, hence eccentricities, hence the centre
set: \(f(\mathrm{centre}(T))=\mathrm{centre}(T')\).  \emph{(<=)} If
\(\Phi(T)=\Phi(T')\), then in the one-centre case both trees rooted at their
centres share a root label, so \(\thmref{c22-ahu}\) supplies a rooted isomorphism,
which is an unrooted one.  In the two-centre case the equal sorted pair means the
endpoint-rooted forms agree in one of the two orientations, and that orientation
is a rooted isomorphism again from \(\thmref{c22-ahu}\).  \emph{(=>)} An
isomorphism \(f\) maps the centre edge to the centre edge, so either
\((T,a)\cong(T',f(a))\) or \((T,a)\cong(T',f(b))\); in the first case
\(x=x'\) and \(y=y'\), in the second \(x=y'\) and \(y=x'\) -- both give the same
sorted pair, and in the one-centre case \(f(c)=c\) forces \(\Phi=\Phi'\).
\end{theorem}
\begin{proofbox}
Only the "centre is preserved" step needs care, so spell it out.  Let
\(\mathrm{ecc}(v)=\max_u \dist(v,u)\) and \(r=\min_v \mathrm{ecc}(v)\).  An
isomorphism satisfies \(\dist(f(v),f(u))=\dist(v,u)\) for all \(u,v\), hence
\(\mathrm{ecc}(f(v))=\mathrm{ecc}(v)\).  So \(f\) maps the set
\(\{v:\mathrm{ecc}(v)=r\}\) onto the corresponding set of \(T'\): a single centre
vertex is mapped to a single centre vertex (it is the \emph{unique} minimizer, so
its image is forced), and a centre edge's two endpoints -- the unique pair
adjacent to each other where both minimize -- are mapped to the other tree's
centre-edge endpoints, in some order.  Nothing else was used, so the two-case
definition of \(\Phi\) is well posed.  Cost: two BFS sweeps for the diameter and
centre (Ch.~12, \chref{c09}), plus at most two labellings per tree at
\(\bigO{n\log n}\); each extra orientation is another \(\bigO{n}\) once the table
is warm. \(\Box\)
\end{proofbox}

\begin{insightbox}[Label-free choices are the only safe ones]
Any rule that mentions a vertex number ("root at \(1\)", "root at the smallest
label", "root at the highest-degree vertex") can be broken by renaming -- \(R\)
versus \(R'\) is that counterexample in miniature.  The centre survives because
it is defined by a distance property invariant under relabelling.  Keep the
principle: \textbf{to canonicalize an unrooted structure, root it at something
the structure itself defines} (centre; or, in other chapters, the diameter's
midpoint).  When the object so defined is an edge rather than a vertex, you get
two candidate roots and must compare a \emph{pair}.
\end{insightbox}

\begin{lstlisting}[caption={Unrooted isomorphism: centre via the diameter (Ch.12 black box), then one shared interned labeller; two centres mean two rooted forms, compared as a sorted pair. ASCII only, identifiers without underscores}, label={lst:c22-unrooted}]
#include <bits/stdc++.h>
using namespace std;
const int MAXN=200005;
int n, nextId;
vector<int> adj[MAXN], kids[MAXN];
map<vector<int>,int> tid;              // ONE table shared by both trees AND by
                                       // both rootings of the same tree
int cl[MAXN], par[MAXN];
int farthest(int src, vector<int>& dist){      // BFS, Ch.12 style
    dist.assign(n+1,-1); queue<int> q; dist[src]=0; q.push(src); int best=src;
    while(!q.empty()){ int u=q.front(); q.pop();
        if(dist[u]>dist[best]) best=u;
        for(int v: adj[u]) if(dist[v]==-1){ dist[v]=dist[u]+1; q.push(v); } }
    return best;
}
void rootat(int rt){                   // iterative BFS rooting (no recursion)
    for(int i=1;i<=n;i++) kids[i].clear();
    vector<int> order; order.push_back(rt); par[rt]=0;
    for(size_t i=0;i<order.size();i++){ int u=order[i];
        for(int v: adj[u]) if(v!=par[u]){ par[v]=u; kids[u].push_back(v); order.push_back(v); } }
    for(int idx=(int)order.size()-1; idx>=0; idx--){
        int u=order[idx]; vector<int> t;
        for(int v: kids[u]) t.push_back(cl[v]);
        sort(t.begin(), t.end());
        map<vector<int>,int>::iterator it=tid.find(t);
        if(it!=tid.end()) cl[u]=it->second; else { cl[u]=nextId++; tid[t]=cl[u]; }
    }
}
// centre: diameter by two BFS sweeps (Ch.12), then walk to the middle of A..B
pair<int,int> centreForm(int start){
    vector<int> d;
    int A=farthest(start,d), B=farthest(A,d);
    int len=d[B];
    rootat(A);                                  // par[] now points toward A
    int u=B; for(int i=0;i<len/2;i++) u=par[u];  // middle vertex of the diameter
    if(len%2==0){ rootat(u); return make_pair(cl[u],cl[u]); }   // one centre
    int c1=u, c2=par[u];                                        // centre EDGE
    rootat(c1); int x=cl[c1];
    rootat(c2); int y=cl[c2];                                    // same tid table
    if(x>y){ int t=x; x=y; y=t; }
    return make_pair(x,y);                       // sorted pair = the canonical form
}
\end{lstlisting}

% ---------------------------------------------------------------
\subsection{Worked verdicts: \((R,R')\), the look-alikes \((T_1,T_2)\), and \(P_6\)}
% ---------------------------------------------------------------

\begin{tracebox}[All three cases, one shared table, ids in interning order]
\textbf{(a) \(R\) vs \(R'\) -- the wrong-root lie, corrected.} \(R\): centre
\(1\), name \(4\) (28.1's run: leaf \(1\), star \(2\), chain-of-two \(3\), root
\(4\)).  \(R'\)'s centre is \(\sigma(1)=2\); labelling \(R'\) at \(2\) against the
same table: children \(5,3,4\); subtree\((5)=\{1,6\}\) two leaves \(\to\) tuple
\((1,1)\to\mathrm{cl}=2\); subtree\((3)=\{7,8\}\to\mathrm{cl}=2\);
subtree\((4)=\{9\}\to\mathrm{cl}=3\); root tuple \((2,2,3)\to\mathrm{cl}=4\) --
\emph{recalled}, not fresh.  \(\Phi(R)=\Phi(R')=4\Rightarrow\boxed{\text{YES}}\),
while the naive rooted-at-\(1\) compare said NO.
\textbf{(b) \((T_1,T_2)\) -- look-alikes, separated.} Both centres are the
\emph{edge} \(\{3,4\}\), so both need two forms. Interning \(T_1\) first:
\(()\to1\) (leaf), \((1)\to2\) (chain-of-two), \((2)\to3\) (chain-of-three),
\((1,2)\to4\).  Rooted at \(3\): children \(2,4,7\) with classes
\(2,3,1\) \(\to\) tuple \((1,2,3)\to\mathrm{cl}=5\).  Rooted at \(4\): children
\(3,5\); subtree\((3)\) has children \(2\) (chain-of-two) and \(7\) (leaf), i.e.\ the
\(\langle\text{leaf},\text{chain}\rangle\) class \(4\), and subtree\((5)\) is a
chain-of-two (class \(2\)) \(\to\) tuple \((2,4)\to6\).
\(\Phi(T_1)=(5,6)\).  Continue on the same table for \(T_2\): the 2-leaf star
\((1,1)\to7\), its one-child parent \((7)\to8\).  Rooted at \(3\): children \(2\)
(a 2-leaf star, class \(7\)) and \(4\) (chain-of-three, class \(3\)) \(\to\) tuple
\((3,7)\to9\).  Rooted at \(4\): children \(3\) (class \(8\)) and \(5\) (class
\(2\)) \(\to\) tuple \((2,8)\to10\).  \(\Phi(T_2)=(9,10)\neq(5,6)\)
\(\Rightarrow\boxed{\text{NO}}\) -- the mechanical version of the distance-profile
argument in figure \(\figref{c22-trap}\).
\textbf{(c) \(P_6\) -- a self-dual centre edge.} Rooted at \(3\): children \(2\)
(chain-of-two) and \(4\) (chain-of-three) \(\to\) tuple \((2,3)\to4\).  Rooted at
\(4\): children \(3\) (chain-of-three) and \(5\) (chain-of-two) \(\to\) the
\emph{same} tuple \((2,3)\to4\).  \(\Phi(P_6)=(4,4)\): when the two halves of the
centre edge are isomorphic the pair degenerates and one labelling suffices --
the shortcut \(P_4,P_6,P_8,\dots\) all enjoy. \(\Box\)
\end{tracebox}

\subsection{Two centres: why ``pick one'' is not a rule}

When the diameter has odd length the centre is an \emph{edge}, and the tree offers
no rule whatsoever for choosing between its two endpoints -- so the canonical form
must keep both.  Forgetting that is not a corner case: it produces a confident NO
on two trees that are the same tree.

\begin{figure}[H]
\centering
\begin{tikzpicture}[diagram, scale=0.92, transform shape,
   tnode/.style={circle, draw=primaryblue!70, fill=primaryblue!8, line width=0.7pt,
                 minimum size=6mm, font=\scriptsize\bfseries, inner sep=1pt},
   warn/.style={circle, draw=warningred, fill=warningred!12, line width=1pt,
                 minimum size=6mm, font=\scriptsize\bfseries, inner sep=1pt}]
  \node[tnode] (q1) at (-3,0) {1};
  \node[tnode] (q2) at (-2,0) {2};
  \node[warn]  (q3) at (-1,0) {3};
  \node[warn]  (q4) at (0,0) {4};
  \node[tnode] (q5) at (1.1,0.9) {5};
  \node[tnode] (q6) at (1.1,-0.9) {6};
  \node[tnode] (q7) at (2.2,-0.9) {7};
  \draw (q1)--(q2)--(q3)--(q4) (q4)--(q5) (q4)--(q6) (q6)--(q7);
  \node[annot, font=\scriptsize, anchor=south, text width=6.6cm, align=center] at (-0.4,1.15)
    {\(Q_1\): chain half on vertex \(3\), branched half on \(4\)};
  \node[tnode] (r1) at (3.4,-4.6) {1};
  \node[tnode] (r2) at (2.4,-4.6) {2};
  \node[warn]  (r4) at (1.4,-4.6) {4};
  \node[warn]  (r3) at (0.4,-4.6) {3};
  \node[tnode] (r5) at (-0.7,-3.7) {5};
  \node[tnode] (r6) at (-0.7,-5.5) {6};
  \node[tnode] (r7) at (-1.8,-5.5) {7};
  \draw (r1)--(r2)--(r4)--(r3) (r3)--(r5) (r3)--(r6) (r6)--(r7);
  \node[annot, font=\scriptsize, anchor=north, text width=6.6cm, align=center] at (1.0,-5.75)
    {\(Q_2\): the same two halves with the \emph{labels} swapped across the centre edge};
\end{tikzpicture}
\caption{The centre-edge trap, with a real counterexample.  \(Q_1\) has edges
\(1\!-\!2,2\!-\!3,3\!-\!4,4\!-\!5,4\!-\!6,6\!-\!7\); \(Q_2\) has
\(1\!-\!2,2\!-\!4,4\!-\!3,3\!-\!5,3\!-\!6,6\!-\!7\).  Both centres are the edge
\(\{3,4\}\) (red), and \(Q_1\cong Q_2\): the bijection that swaps \(3\) and
\(4\) and fixes every other vertex maps \(Q_1\)'s edges onto \(Q_2\)'s one by
one -- \(1\!-\!2\!\mapsto\!1\!-\!2\), \(2\!-\!3\!\mapsto\!2\!-\!4\),
\(3\!-\!4\!\mapsto\!4\!-\!3\), \(4\!-\!5\!\mapsto\!3\!-\!5\),
\(4\!-\!6\!\mapsto\!3\!-\!6\), \(6\!-\!7\!\mapsto\!6\!-\!7\) -- exactly the listed
set of \(Q_2\).  The tempting rule ``root both at the smaller-labelled centre endpoint
\(3\)'' fails: \(\Phi_3(Q_1)=\langle\)chain-of-two,
\(\langle\text{leaf},\text{chain-of-two}\rangle\rangle\) has \emph{two} children
while \(\Phi_3(Q_2)=\langle\)leaf, chain-of-two, chain-of-three\(\rangle\) has
\emph{three} -- different names, so a single-endpoint rule answers NO on a YES
instance.  The sorted pair of \(\defref{c22-canroot}\) is the fix:
\(\{\Phi_3,\Phi_4\}\) is the same multiset for both trees, only with the roles of
\(3\) and \(4\) exchanged.}
\label{fig:c22-pair}
\end{figure}

% ---------------------------------------------------------------
\subsection{Problem Lab: CF 1800G Symmetree -- can the children be paired into mirrors?}
% ---------------------------------------------------------------

\solvedproblem{c22-symmetree}{Symmetree}{CF 1800G}
  {https://codeforces.com/problemset/problem/1800/G}{CF 2200 L4}
  {one AHU pass gives every child's shape name; a tree is symmetrical iff at every vertex at most one child \emph{class} occurs an odd number of times (\(\thmref{c22-sym}\)); iterate the rootify so \(t\cdot n\) fits}

\begin{problemstatement}[Symmetree (statement, constraints and all six samples verified against the live CF page)]
The tree has \(n\) vertices, \(1\le n\le 2\cdot10^{5}\), rooted at vertex \(1\);
\(t\) test cases (\(1\le t\le10^{4}\)) with \(\sum n\le 2\cdot10^{5}\), each
followed by \(n-1\) undirected edges.  Kid wants the tree to look
\emph{symmetrical}: is there \textbf{an order of the children of every vertex}
such that the leftmost child's subtree is the mirror image of the rightmost
child's, the second-left mirrors the second-right, and so on -- and if a vertex
has an odd number of children, the middle one's subtree must itself be
symmetrical?  Print YES or NO per test case.  (A mirror image of a rooted tree is
just its left-right reversal, so "subtree \(A\) mirrors subtree \(B\)'' \emph{is}
"rooted-isomorphic'', \(\defref{c22-iso}\).)
\end{problemstatement}

\begin{figure}[H]
\centering
\begin{tikzpicture}[diagram, scale=0.95, transform shape,
   tnode/.style={circle, draw=primaryblue!70, fill=primaryblue!8, line width=0.7pt,
                 minimum size=6.4mm, font=\scriptsize\bfseries, inner sep=1pt},
   sub/.style={circle, draw=successgreen!70, fill=successgreen!14, line width=1pt,
                 minimum size=6.4mm, font=\scriptsize\bfseries, inner sep=1pt}]
  \node[tnode] (c1) at (0,0) {1};
  \node[sub]   (c5) at (-3,-1.4) {5};
  \node[tnode] (c2) at (0,-1.4) {2};
  \node[sub]   (c6) at (3,-1.4) {6};
  \node[tnode] (c3) at (-1.1,-2.8) {3};
  \node[tnode] (c4) at (1.1,-2.8) {4};
  \draw (c1)--(c5) (c1)--(c2) (c1)--(c6) (c2)--(c3) (c2)--(c4);
  \node[annot, font=\scriptsize, anchor=south, text width=7.2cm, align=center] at (0,0.7)
    {pair \(5\leftrightarrow6\) (both leaves), middle \(2\) -- which is itself symmetrical};
\end{tikzpicture}
\caption{Sample 1 of CF 1800G (\(n=6\), edges \(1\!-\!5,1\!-\!6,1\!-\!2,2\!-\!3,
2\!-\!4\)): the answer is YES.  The root's three children split into one mirrored
pair (green) and one middle child, whose own subtree must be symmetrical -- the
recursion that \(\thmref{c22-sym}\) turns into a parity count.}
\label{fig:c22-symmetree}
\end{figure}

\begin{attemptgate}[Attempt it yourself first -- \lv{4}, give yourself 30 minutes]
What exactly must two children that play the mirror roles have in common, and what
happens to the child left in the middle?
\peek{"Mirror image" of a rooted tree = left-right reversal.  Reversal does not
change the unordered structure, so \(A\) mirrors \(B\) iff \(A\) and \(B\) are
rooted-isomorphic -- you already own the key that makes this an integer compare
(\(\secref{c22-ahu}\)).}
\nudge{At each vertex, group the children by their shape name and look at the
parity of each group's size.  Pairs consume children two at a time, so every
group must have even size... except possibly one group, whose leftover child
sits in the middle.}
\fullhint{Run one postorder AHU labelling with a table shared inside the test
case, and at the same visit count how many child-classes occur an odd number of
times.  Claim: the tree is symmetrical iff \textbf{every} vertex has \(\le 1\)
odd class (\(\thmref{c22-sym}\)).  Sufficiency is induction on height: pair the
children class-by-class, and the single unpaired middle child has class \(C\);
every vertex of its subtree still obeys the rule, so by induction that subtree is
symmetrical and may stand in the middle.  Necessity: a legal ordering pairs equal
classes and leaves at most one middle, so only one class can be odd.  Iterate the
rooting (BFS order) -- \(\sum n=2\cdot10^{5}\) and a chain will otherwise blow the
stack; clear \(\mathit{tid}\) per test case, since ids are only meaningful inside
one case.}
\end{attemptgate}

\solstep{Understanding}
Two claims do all the work.  First, "mirror image" is isomorphism: reversing the
child order of a rooted tree yields the same unordered structure, so
\(A\) mirrors \(B\) iff \(\mathrm{cl}(A)=\mathrm{cl}(B)\).  Second, the definition
is \emph{local once shapes are known}: the only way a vertex can fail is by having
two children that cannot be paired and cannot both be middles -- i.e.\ two odd
shape classes.  The global "does an ordering exist" question collapses into a
per-vertex parity count, which costs one extra scan of a sorted array you already
built.

\solstep{Key Observations}
\begin{observation}[Mirror \(=\) isomorphic]
For rooted trees, \(\mathrm{mirror}(A)\cong A\) and "\(A\) mirrors \(B\)''
\(\iff\) \((A,r_A)\cong(B,r_B)\) (\(\defref{c22-iso}\)).  So the pairing key is
the canonical label \(\mathrm{cl}\) -- or its hash
(\(\defref{c22-shapehash}\)) -- and nothing about left/right has to be
represented at all.
\end{observation}
\begin{observation}[Parity form of the symmetry test]
A vertex is \emph{pairable} iff at most one of its child shape-classes has odd
multiplicity.  The tree is symmetrical iff every vertex is pairable AND, whenever
a class is odd, that class is itself symmetrical -- which, because symmetry is a
property of the class (isomorphic subtrees share it), follows inductively from
every vertex being pairable.  Hence: symmetrical \(\iff\) all vertices pairable.
\end{observation}

\begin{theorem}{Symmetry is one odd class per vertex}{c22-sym}
A rooted tree \(T\) is symmetrical in the sense of CF 1800G iff for every vertex
\(v\), at most one shape-class among \(\children(v)\) occurs an odd number of
times.
\end{theorem}
\begin{proofbox}
\emph{(=>)} A legal ordering pairs children of equal class and leaves at most one
child (the middle) unpaired, so all classes but possibly one are consumed two at
a time.  \emph{(<=)} Induct on the height.  At \(v\), pair the children
class-by-class -- equal class means rooted-isomorphic, and paired subtrees have
their own legal orderings by induction, so mirror one of each pair against the
other.  If one class has odd size, put a child \(m\) of that class in the middle:
every vertex of \(\subtree(m)\) still obeys the parity rule, so \(\subtree(m)\) is
symmetrical by induction, exactly as the definition demands. \(\Box\)
\end{proofbox}

\solstep{The Approach Ladder}
\begin{center}\renewcommand{\arraystretch}{1.25}\setlength{\tabcolsep}{3pt}
\begin{tabular}{@{}c p{0.26\linewidth} p{0.30\linewidth} l c@{}}
\toprule
\textbf{\#} & \textbf{Approach} & \textbf{Idea} & \textbf{Time} & \textbf{Verdict}\\
\midrule
1 & \lv{2}\ try every child order & permute children, test the picture & \(\bigO{n!\text{--}2^n}\) & \tle \\
2 & \lv{3}\ pair by subtree size & equal size can stand in for equal shape & \(\bigO{n\log n}\) & \textbf{WA} \\
3 & \lv{4}\ pair by AHU class at the root only & check the root, assume the middle works & \(\bigO{n\log n}\) & \textbf{WA} \\
4 & \lv{4}\ \textbf{pair by class at every vertex} & per-vertex odd-class count \(\le1\) & \(\bigO{n\log n}\) & \ac \\
\bottomrule
\end{tabular}\end{center}
Rung~2 is not a shape test at all: under it a \(3\)-node chain and a \(3\)-node
star are interchangeable children, so it accepts orderings that cannot be drawn.
Rung~3 forgets that the middle child is not free -- it inherits the whole
obligation, which is where the official samples 4 and 5 fail, one level below a
root that looks perfectly pairable.

\solstep{Solution}
Per test case: root at \(1\) iteratively (BFS order); walk the order backwards; at
each vertex collect its children's \(\mathrm{cl}\) values, sort them, count the
runs of odd length, and flag the vertex when that count exceeds one; then intern
the sorted tuple to obtain \(\mathrm{cl}(u)\).  Print YES iff no vertex was ever
flagged.

\solstep{Implementation}
\begin{lstlisting}[caption={CF 1800G. One shared interned labeller + a per-vertex odd-run count. Iterative rooting (sum n = 2e5). ASCII only, identifiers without underscores}, label={lst:c22-symmetree}]
#include <bits/stdc++.h>
using namespace std;
const int MAXN=200005;
int n, cl[MAXN];
vector<int> adj[MAXN], kids[MAXN];
map<vector<int>,int> tid;
int main(){
  int tc; scanf("%d",&tc);
  while(tc--){
    scanf("%d",&n);
    for(int i=1;i<=n;i++){ adj[i].clear(); kids[i].clear(); }
    tid.clear();
    for(int i=1;i<n;i++){ int u,v; scanf("%d%d",&u,&v); adj[u].push_back(v); adj[v].push_back(u); }
    vector<int> order; order.push_back(1);
    vector<int> seen(n+1,0); seen[1]=1;
    for(size_t i=0;i<order.size();i++){                 // BFS/DFS order: parent before child
      int u=order[i];
      for(int v: adj[u]) if(!seen[v]){ seen[v]=1; kids[u].push_back(v); order.push_back(v); }
    }
    bool ok=true;
    for(int idx=(int)order.size()-1; idx>=0; idx--){    // children before parent
      int u=order[idx]; vector<int> t;
      for(int v: kids[u]) t.push_back(cl[v]);
      sort(t.begin(), t.end());
      int odd=0;
      for(size_t j=0;j<t.size();){
        size_t k=j; while(k<t.size() && t[k]==t[j]) k++;
        if( (k-j) & 1 ) odd++;
        j=k;
      }
      if(odd>1) ok=false;                               // two children cannot be paired
      map<vector<int>,int>::iterator it=tid.find(t);
      if(it!=tid.end()) cl[u]=it->second;
      else { int id=(int)tid.size()+1; tid[t]=id; cl[u]=id; }
    }
    puts(ok?"YES":"NO");
  }
  return 0;
}
\end{lstlisting}

\solstep{Trace}
\begin{tracebox}[The six official samples, each with the reason it says what it says]
\textbf{1.} \(n{=}6\), \(1\!-\!5,1\!-\!6,1\!-\!2,2\!-\!3,2\!-\!4\).  Root \(1\):
children \(5,6\) (leaf class, run \(2\) -- even) and \(2\) (star class, run \(1\)
-- odd) \(\Rightarrow\) one odd class, fine; vertex \(2\): two leaves, run \(2\),
zero odd \(\Rightarrow\) \(\mathbf{YES}\).
\textbf{2.} \(n{=}7\), \(1\!-\!5,1\!-\!3,3\!-\!6,1\!-\!4,4\!-\!7,4\!-\!2\).  Root
\(1\)'s children: \(5\) (leaf), \(3\) (chain-of-two), \(4\) (star of two leaves)
-- three classes, three odd runs \(>1\) \(\Rightarrow\Rightarrow\mathbf{NO}\).
\textbf{3.} \(n{=}9\), \(1\!-\!2,2\!-\!4,2\!-\!3,3\!-\!5,1\!-\!7,7\!-\!6,7\!-\!8,8\!-\!9\).
subtree\((2)\): children \(4\) (leaf) and \(3\) (chain-of-two) \(\to\) class
\(\langle\text{leaf},\text{chain}\rangle\); subtree\((7)\): children \(6\) (leaf)
and \(8\) (chain-of-two) \(\to\) \emph{same} class.  Root has those two children
only, one run of length \(2\) -- zero odd \(\Rightarrow\mathbf{YES}\).
\textbf{4.} \(n{=}10\), \(2\!-\!9,9\!-\!10,2\!-\!3,6\!-\!7,4\!-\!3,1\!-\!2,3\!-\!8,2\!-\!5,6\!-\!5\).
Root \(1\): single child \(2\) (one odd run, passes).  Vertex \(2\)'s children:
\(9\) (chain-of-two, since \(9\!-\!10\)), \(3\) (star: \(3\!-\!4,3\!-\!8\)), \(5\)
(chain-of-three: \(5\!-\!6\!-\!7\)) -- three distinct classes, three odd runs
\(\Rightarrow\mathbf{NO}\) (rung~3 would have stopped at the happy root).
\textbf{5.} \(n{=}10\), \(3\!-\!2,8\!-\!10,9\!-\!7,4\!-\!2,8\!-\!2,2\!-\!1,4\!-\!5,6\!-\!5,5\!-\!7\).
Vertex \(2\)'s children: \(3\) (leaf), \(8\) (chain-of-two: \(8\!-\!10\)), \(4\)
(\(\langle\)chain-of-two,\(\star\rangle\) via \(5\) with \(6\) and \(7\!-\!9\)) --
three distinct classes again \(\Rightarrow\mathbf{NO}\).
\textbf{6.} \(n{=}1\): the root has zero children, zero odd runs
\(\Rightarrow\mathbf{YES}\).  Print \(\mathbf{YES\ NO\ YES\ NO\ NO\ YES}\) -- the
official output. \(\Box\)
\end{tracebox}

\begin{warningbox}[Three ways 1800G goes wrong]
\begin{enumerate}[nosep,leftmargin=1.6em]
  \item \textbf{Pairing by size (rung 2).} \(\size\) is not a shape name: the
        chain and the star on four vertices tie on size and lose on shape.  If
        your key is not \(\mathrm{cl}\) (or a hash of shape,
        \(\defref{c22-shapehash}\)) the verdict is a coin flip.
  \item \textbf{Checking only the root (rung 3).} The middle child is not free --
        it inherits the whole obligation.  Samples 4 and 5 fail one level down,
        which is exactly what "at most one odd class at \emph{every} vertex"
        buys you.
  \item \textbf{Recursion and stale tables.} \(\sum n=2\cdot10^{5}\) on a chain
        overflows the default stack, and reusing \(\mathit{tid}\) across test
        cases with a stale \(\mathit{nextId}\) silently merges classes from
        different cases -- clear the table per case (ids are only comparable
        inside one table).
\end{enumerate}
\end{warningbox}

\begin{hardeningbox}
\begin{itemize}[nosep,leftmargin=1.4em]
  \item \textbf{Hash instead of intern, if you prefer.} Replace \(\mathrm{cl}\) by
        the commutative depth-seeded hash \(h\) of \(\defref{c22-shapehash}\): the
        parity test reads the same, sorting \(64\)-bit values, and no table is
        cleared per case.  You inherit the collision ledger
        (\(\thmref{c22-collision}\)) -- for a YES/NO judge problem with
        \(n\le2\cdot10^{5}\) a single \(64\)-bit hash is the usual contest
        choice; the deterministic version above is what you want if the answer
        must be provable.
  \item \textbf{The same test finds the axis.} The odd class at each vertex
        \emph{names} the middle child; walking those pointers from the root gives
        the symmetry axis, which is what CF 1182/D ``Complete Mirror'' asks you to
        print -- try it after this section (it is solved in the advanced tree-DP
        chapter, so treat 1800G as its warm-up).
  \item \textbf{Odd-Run counting without a map.} Sorting the child labels is
        enough for the parity test -- if you never need the ids again, drop the
        interning entirely and keep only per-level sorted arrays (still
        \(\bigO{n\log n}\)).
\end{itemize}
\end{hardeningbox}

\begin{recallbox}[28.3 in one breath]
\emph{Root at the thing the tree defines: the centre is fixed by a distance
property, so it survives every renaming; one centre \(\to\) compare
\(\mathrm{cl}\) at it, centre edge \(\to\) compare the \textbf{sorted pair} of the
two endpoint forms (\(\thmref{c22-centreiso}\), \(\coderef{c22-unrooted}\)).  \(R\)
against \(R'\): rooted-at-\(1\) NO, centre-rooted YES (both name \(4\));
\((T_1,T_2)\): both centre-edged at \(\{3,4\}\) yet pairs \((5,6)\) against
\((9,10)\) -- NO with certainty; \(P_6\): pair \((4,4)\), the self-dual shortcut.
Mirror-symmetry (CF 1800G) is the same pairing instinct one level down: children
pair by equal shape class, at most one odd class per vertex, middle child
included (\(\thmref{c22-sym}\), \(\coderef{c22-symmetree}\)) -- six samples,
YES NO YES NO NO YES.}
\selftest{
  \item Why can an isomorphism never move the centre? \emph{(It preserves distances, hence eccentricities, hence the set of minimizers -- the centre.)}
  \item Centre is the edge \(\{a,b\}\): what exactly is compared? \emph{(The sorted pair of the two endpoint-rooted canonical names; equality of pairs iff isomorphic.)}
  \item Root passes the parity test but the answer is NO -- how? \emph{(The middle child's own subtree fails one level down; the condition must hold at every vertex. Samples 4/5 of 1800G.)}
}
\end{recallbox}
